\begin{apartats}

\begin{tria}{limits-infinit-tipus}
\itemtria{un-limit}{0,75}{1,25}
Calcula el límit següent:
\[
\lim_{x\to+\infty}\frac{3x^2+1}{x-4}
\]

\begin{solucio}
El grau del numerador ($2$) és més gran que el del denominador ($1$). El quocient dels termes de
grau màxim, $\dfrac{3x^2}{x}=3x$, tendeix a $+\infty$: el límit és $+\infty$.
\end{solucio}

\itemtria{dos-tipus}{1,25}
Calcula els límits següents:
\[
\lim_{x\to+\infty}\frac{3x^2+1}{x-4}
\qquad
\lim_{x\to+\infty}\frac{3x+1}{x-4}
\]

\begin{solucio}
i) El grau del numerador ($2$) és més gran que el del denominador ($1$). El quocient dels termes de
grau màxim, $\dfrac{3x^2}{x}=3x$, tendeix a $+\infty$: el límit és $+\infty$.\\
ii) Numerador i denominador tenen el mateix grau ($1$). El límit val el quocient dels
coeficients de grau màxim: $\dfrac{3}{1}=3$.
\end{solucio}

\itemtria{quatre-tipus}{2}
Calcula els límits següents:
\[
\lim_{x\to+\infty}\frac{3x^2+1}{x-4}
\qquad
\lim_{x\to+\infty}\frac{3x+1}{x-4}
\qquad
\lim_{x\to-\infty}\frac{x^3+1}{1-x}
\qquad
\lim_{x\to-\infty}\frac{11}{x-2}
\]

\begin{solucio}
i) El grau del numerador ($2$) és més gran que el del denominador ($1$). El quocient dels termes de
grau màxim, $\dfrac{3x^2}{x}=3x$, tendeix a $+\infty$: el límit és $+\infty$.\\
ii) Numerador i denominador tenen el mateix grau ($1$). El límit val el quocient dels
coeficients de grau màxim: $\dfrac{3}{1}=3$.\\
iii) El grau del numerador ($3$) és més gran que el del denominador ($1$). El quocient dels termes
de grau màxim, $\dfrac{x^3}{-x}=-x^2$, tendeix a $-\infty$ quan $x\to-\infty$: el límit és $-\infty$.\\
iv) El grau del denominador ($1$) és més gran que el del numerador ($0$): el límit val $0$.
\end{solucio}
\end{tria}

\begin{nomesllarg}
\apartat{0,75}
Calcula el límit següent:
\[
\lim_{x\to-\infty}\frac{1-x^4}{2x^4-x^2+5}
\]

\begin{solucio}
Numerador i denominador tenen el mateix grau ($4$): el límit és el quocient dels coeficients
principals, $\dfrac{-1}{2}=-\dfrac12$.
\end{solucio}
\end{nomesllarg}

\begin{tria}{determina-a}
\itemtria{amb-reflexio}{1}{1,25}
Considera la funció
\[
f(x)=\frac{ax^2-2x}{4x^2+1},
\]
on $a$ és un nombre real. Determina $a$ perquè
\[
\lim_{x\to+\infty}f(x)=\frac12 .
\]
Existeix algun valor de $a$ per al qual aquest límit valgui $0$? Justifica-ho.

\begin{solucio}
Si $a\neq0$, numerador i denominador tenen grau $2$ i el límit val $\dfrac{a}{4}$.
Imposant $\dfrac{a}{4}=\dfrac12$ s'obté $\boxed{a=2}$.\\
Si $a=0$, el numerador és $-2x$, de grau $1$, i el límit val $0$. Per tant, \textbf{sí}: el
límit val $0$ quan $a=0$, i només en aquest cas, perquè si $a\neq0$ val $\dfrac a4\neq0$.
\end{solucio}

\itemtria{sense-reflexio}{0,5}
Considera la funció
\[
f(x)=\frac{ax^2-2x}{4x^2+1},
\]
on $a$ és un nombre real. Determina $a$ perquè
\[
\lim_{x\to+\infty}f(x)=3 .
\]

\begin{solucio}
Numerador i denominador tenen grau $2$ i el límit val $\dfrac{a}{4}$.
Imposant $\dfrac{a}{4}=3$ s'obté $\boxed{a=12}$.
\end{solucio}
\end{tria}

\end{apartats}
